% O7. Orders that share a posterior.
% First paragraph: Appendix A.2, in place of appendix/A2_beliefs_trading.tex:165-170 (the paragraph on constructed
% tuples), after the proof of Lemma A.2 and before Lemma A.3.
% source: scan/reduction_A3_A5.md, part 2 ("At a shared posterior, prices agree exactly when entry probabilities
% agree") and part 4, item 1 (the equal-entry qualification under the second challenger rule); referee report W3,
% "Shared Candidate 1", section 4 (pooled sale and zero order); PLAN.md Part 5, O7.
% Second paragraph: Appendix A.5, in the verification of recruiting research, at the condition (E2).

\paragraph{Constructed tuples.}
Let \(\DT>0\). Take a profile, the posteriors \(\mu(x)\) that (E1) gives it, and entry \(e(x)\in[\rho,1]\) after
each order that is nondecreasing in the posterior: \(e(x)\ge e(x')\) whenever \(\mu(x)>\mu(x')\). Set
\(P(x)=t_0+e(x)\,w(\mu(x))\). Then \(P(x)>P(x')\) whenever \(\mu(x)>\mu(x')\), and two orders have the same price if
and only if they have the same posterior and the same entry. Hence \(\zeta(P(x))=[e(x)-\rho]/(1-\rho)\) is a
preparation rule on the image of \(P\), with values in \([0,1]\), and the belief that (E2) assigns to the price
\(P(x)\) is \(\mu(x)\). Let also \(g_H>g_L\). When an indifferent high-cost challenger prepares, entry that meets
Lemma~\ref{lem:prices}(d) is a function of the posterior, so orders with the same posterior have the same price.
When an indifferent high-cost challenger may prepare with any probability, orders with the posterior \(\tau\)
and different entry have different prices, each revealing \(\tau\).

\begin{proof}
Let \(\mu(x)>\mu(x')\). Then \(e(x)\ge e(x')\), and \eqref{eq:a2-increment} gives
\(P(x)-P(x')\ge e(x)\DT[\mu(x)-\mu(x')]\ge\rho\DT[\mu(x)-\mu(x')]>0\), because \(w(\mu)\ge w_L\ge0\). Let
\(\mu(x)=\mu(x')\). Then \(P(x)-P(x')=[e(x)-e(x')]\,w(\mu(x))\), and \(w(\mu(x))=w_L+\DT\mu(x)>0\), because
\(\mu(x)>0\) by Lemma~\ref{lem:beliefs}(a). Equal prices therefore force equal posteriors and equal entry, and the
price formula gives the converse. So \(\zeta\) is a function of the price, and \(e(x)\in[\rho,1]\) puts its values in
\([0,1]\). At a price \(\bar P\), all orders with the price \(\bar P\) share one posterior, so by \eqref{eq:odds}
they share one likelihood ratio \(f_H(x)/f_L(x)\). The ratio of sums in the step on beliefs at a price, in the
proof of Lemma~\ref{lem:beliefs}, equals that likelihood ratio, and (E2) gives the common posterior. Under the
first challenger rule, Lemma~\ref{lem:prices}(d) fixes entry at every posterior, \(\tau\) included. Under the second,
the identity at equal posteriors gives distinct prices to distinct entries at \(\tau\), and every order with one of
these prices has the posterior \(\tau\).

\end{proof}

Each construction below applies this paragraph to its stated entry to obtain its price beliefs. Where two orders
of a stated tuple share a posterior, the tuple gives them the same entry, so they share a price, and the belief at
that price is their common posterior.

\paragraph{The sale and the zero order in recruiting research.}
At the prior \(1/2\), when an indifferent investor may mix, let \(\tau>1/2\), and take recruiting research with \(0<\hat\lambda<1\), in which class
\(H\) buys after a success and class \(L\) sells with probability \(\sigma_L(-1)\) and abstains otherwise.
The sale and the zero order have entry \(\rho\) under either challenger tie rule. They share a posterior,
and hence a price, exactly when
\[
\sigma_L(-1)=\frac{\beta}{1-\beta}.
\]
In the example, during recruiting research at the cap, this happens at exactly one strength,
\[
r=\fr\!\left(\frac{k}{\rho}\left[2+\frac{(\eu-1)\beta}{1-\beta}\right]\right)
\approx\reg{r_pool_B}.
\]
The common posterior is \(\reg{pool_mu_B}\), and the common price is \(\reg{pool_price_B}\).

\begin{proof}
With \(\hat b=(1-\hat\lambda)\beta\) and \(\hat b_0=(1-\hat\lambda)(1-2\beta)\), Bayes' rule gives
\[
\mu(-1)=\frac{\hat b}{2\hat b+\hat\lambda\sigma_L(-1)}\le\frac12,\qquad
\mu(0)=\frac{\hat b_0}{2\hat b_0+\hat\lambda[1-\sigma_L(-1)]}\le\frac12.
\]
Both posteriors lie strictly below \(\tau\), so both orders have entry \(\rho\) with no challenger tie. A purchase
fully recruits. Entry is nondecreasing in the posterior, and the paragraph on constructed tuples gives the
price beliefs. Cross multiplication gives equal nonpurchase posteriors exactly when
\(\hat b[1-\sigma_L(-1)]=\hat b_0\sigma_L(-1)\). Since
\(\hat b/(\hat b+\hat b_0)=\beta/(1-\beta)\), this is the stated sale probability, independently of research.

On the example's capped segment where the sale is mixed, the sale earns zero, so
\[
\hat s_-(r)=\bar b\left[\frac{\rho\DT(r)}{k}-2\right],\qquad
\sigma_L(-1)=\frac{\hat s_-(r)}{\lb},\qquad
\frac{\lb}{\bar b}=\eu-1.
\]
The spread rises strictly, so the sale probability rises strictly from zero to one as \(r\) ranges from
\(r_N\) to \(r_S\). The number \(\beta/(1-\beta)\) lies strictly between zero and one. It is attained once,
at the spread \((k/\rho)[2+(\eu-1)\beta/(1-\beta)]\). Inverting
\(\DT(r)=(r-\ell)^2/(2r)\) gives the stated strength. The example's ordering
\(r_{\lb}<r_N<r_S<r_C\) puts this strength on the capped recruiting segment.
The sale's zero return gives \(\mu(-1)=k/[\rho\DT(r)]\), and its price is
\(t_0+\rho(t_L-t_0)+k\); these give the stated common posterior and price.
\end{proof}
